\section{The Admission Grammar}
\label{sec:grammar}

The grammar names what a substrate is willing to admit. We describe it informally and cite the underlying formalisation in companion work~\cite{programmableSovereignty2026}. The point is to give the reader a clear picture of the typed object that admission operates on, without requiring acceptance of the companion paper as a premise. This workshop paper does not prove a new reversible-action theorem; its claims about rollback and executive-action commitments are grammar commitments in this application layer, bounded by the cited parent substrate and by the assumptions in Section~\ref{sec:threat}.

\paragraph{Action classes.}
A tool call belongs to one of three action classes. \emph{Observational} calls read state without mutating it (a database SELECT, a file read, an idempotent HTTP GET). \emph{Reversible} calls mutate state with a known inverse executor: a file quarantine has a release, a row insert has a delete, a process suspend has a resume. \emph{Destructive} calls mutate state without a typed inverse: a permanent record deletion that empties a backup, a transaction broadcast to a public ledger, an email sent to an external recipient.

The classification is action-dependent rather than agent-dependent, and is the substrate's commitment, not the agent's~\cite{programmableSovereignty2026}. The substrate decides, for each tool a server exposes, which class each method falls into. An unclassified method is destructive by default (fail-closed).

\paragraph{The admission envelope.}
The substrate accepts a tool call wrapped in an admission envelope that carries four fields beyond the call's intrinsic arguments: (i) an authorisation chain (the capability token under which the agent acts), (ii) an action class (observational, reversible, or destructive), (iii) a positive TTL by construction (the maximum duration the call's effects may stand before the substrate considers them stale), and (iv) for reversible and destructive classes, a rollback receipt slot (an opaque record reserved at admission to be filled at rollback time).

The envelope is typed. A reversible call without a TTL is not constructable: the substrate's input parser rejects it before any executor sees it. A destructive call without both a TTL and a bilateral cosignature is not constructable. The envelope is the syntactic constraint, not a runtime policy. The bilateral cosignature predicate for irreversible variants is developed concretely in a companion construction (anonymized for review): a two-party attestation whose admission predicate composes a small set of gates over canonical bytes.

\begin{figure}[t]
\centering
\resizebox{\textwidth}{!}{%
\begin{tikzpicture}[
  font=\footnotesize,
  envelope/.style={draw, thick, rounded corners=2pt, fill=gray!8, align=center, inner sep=4pt, minimum width=18mm, minimum height=12mm},
  gate/.style={draw, thick, diamond, aspect=2.4, fill=blue!7, draw=blue!55!black, align=center, inner sep=1pt, minimum width=26mm, minimum height=14mm, font=\scriptsize},
  accept/.style={draw, thick, rounded corners=2pt, fill=green!10, draw=green!55!black, align=center, inner sep=4pt, minimum width=20mm, minimum height=12mm},
  reject/.style={draw, thick, rounded corners=2pt, fill=red!7, draw=red!60!black, align=center, inner sep=3pt, minimum width=22mm, minimum height=8mm, font=\scriptsize},
  pass/.style={-{Latex[length=2mm]}, thick},
  fail/.style={-{Latex[length=2mm]}, thick, red!55!black, dashed},
  passlabel/.style={font=\scriptsize, above, inner sep=1pt},
  faillabel/.style={font=\scriptsize\itshape, red!55!black, right=1mm, inner sep=1pt}
]

\node[envelope] (env) at (0, 0) {Envelope\\emitted};
\node[gate, right=10mm of env] (g1) {Gate 1\\capability};
\node[gate, right=10mm of g1] (g2) {Gate 2\\class};
\node[gate, right=10mm of g2] (g3) {Gate 3\\TTL $> 0$};
\node[gate, right=10mm of g3] (g4) {Gate 4\\rollback slot};
\node[accept, right=10mm of g4] (exec) {Executed\\+ receipt};

\draw[pass] (env) -- (g1);
\draw[pass] (g1) -- node[passlabel]{pass} (g2);
\draw[pass] (g2) -- node[passlabel]{pass} (g3);
\draw[pass] (g3) -- node[passlabel]{pass} (g4);
\draw[pass] (g4) -- node[passlabel]{pass} (exec);

\node[reject, below=16mm of g1] (r1) {\texttt{cap-denied}};
\node[reject, below=16mm of g2] (r2) {\texttt{class-mismatch}};
\node[reject, below=16mm of g3] (r3) {\texttt{ttl-invalid}};
\node[reject, below=16mm of g4] (r4) {\texttt{rollback-missing}};

\draw[fail] (g1.south) -- node[faillabel]{fail} (r1.north);
\draw[fail] (g2.south) -- node[faillabel]{fail} (r2.north);
\draw[fail] (g3.south) -- node[faillabel]{fail} (r3.north);
\draw[fail] (g4.south) -- node[faillabel]{fail} (r4.north);

\node[draw=red!55!black, dashed, thick, rounded corners=2pt, fit=(r1) (r2) (r3) (r4), inner sep=4pt, label={[font=\scriptsize\itshape, red!55!black]below:rejection receipt with named code}] {};

\end{tikzpicture}%
}
\caption{Tool-call lifecycle with substrate-side admission gates. The model emits a candidate envelope; the substrate evaluates each gate left-to-right; rejection at any gate emits a named code; acceptance produces an admission receipt that closes the action chain. Gate 4 carries the rollback-receipt-slot check for reversible classes and the bilateral-cosignature check for destructive classes.}
\label{fig:tool-call-lifecycle}
\end{figure}

\paragraph{The four operations.}
Admission proceeds by four decidable operations (Figure~\ref{fig:tool-call-lifecycle}). (1) \emph{Capability check}: the authorisation chain attenuates a valid trust root and the residual permits the call's intended effect on the named subject. (2) \emph{Class check}: the call's classified action class matches the envelope's declared class. (3) \emph{TTL check}: the envelope's declared TTL is positive and bounded by the substrate's per-class maximum. (4) \emph{Rollback slot}: for reversible and destructive classes, the envelope carries a typed rollback receipt slot; for destructive classes, a bilateral cosignature is also present from both the agent's operator polity and the tool server's host polity.

\paragraph{Structural commitments carried by the substrate.}
The grammar uses three structural commitments from the companion substrate~\cite{programmableSovereignty2026}. We state them here as application-layer obligations, not as new machine-checked results in this paper. If the companion remains unpublished at submission time, this section still supplies the local contract the workshop argument needs: admission envelopes are classified, TTL-bound, and routed through bilateral consent at the destructive boundary.

\emph{Commitment 1 (Bounded executive action carries TTL and rollback slot).} A reversible or destructive call's envelope is well-typed only if it carries a positive TTL witness and a rollback receipt slot. The claim is a grammar-level invariant for this workshop construction; its substantive use is the type-checker's refusal to construct envelopes without the discipline.

\emph{Commitment 2 (Rollback admission composes with refinement).} When the admission predicate is amended (the policy changes), a rollback receipt that closes a previously-admitted call remains admissible under the amended predicate, provided the amendment is backward-refining. Rollback obligations do not vanish across policy changes.

\emph{Commitment 3 (Treaty admission iff predicate intersection).} A destructive call admits only if both the operator polity and the host polity admit it under their respective predicates. This is the structural reason a single-party operator manipulation does not suffice to admit a destructive tool call against a counterparty whose policy forbids it.

\paragraph{What the grammar does not claim.}
The grammar does not claim that every tool call should be destructive-class. Observational and reversible classes are the common case and admit on much weaker grounds; the discipline matters at the destructive boundary. The grammar does not claim that the substrate can synthesise a rollback executor for an intrinsically irreversible call. It claims only that, for such a call, admission requires bilateral cosignature in addition to the typed witness, and the cosignature is a public commitment that the call has been admitted as irreversible by both parties.

The central commitment is that a destructive tool call is not constructable in the typed runtime path without its audit slot and destructive-class bilateral cosignature. Misalignment in the issuing model does not produce envelopes that bypass the host polity's independent admission. The two-party structural check, not the operator's diligence, is what the substrate enforces.
